\section{Double Depth First Search}

\subsection{Preliminaries}

\newcommand{\ddfs}{\textsf{\upshape ddfs}}
\newcommand{\digraph}{\DAG}
\newcommand{\level}{\textsf{\upshape level}}
\renewcommand{\choice}{\textsf{\upshape choose}}
\newcommand{\visa}{\ensuremath{\vis_1}}
\newcommand{\visb}{\ensuremath{\vis_2}}
\newcommand{\patha}{\ensuremath{\path_1}}
\newcommand{\pathb}{\ensuremath{\path_2}}
\newcommand{\Plevel}[1]{\ensuremath{P_{#1}}}
\newcommand{\butlast}{\textsf{\upshape butlast}}

A digraph $\digraph$ is a set of ordered pairs of vertices, each of which corresponding to a directed edge.
The set of vertices $\vertices(\DAG)$ is defined to be the set $\bigcup\{\{\vertexa,\vertexb\}\mid(\vertexa,\vertexb)\in\digraph\}$.
A list of vertices $\vertexa\vertexb\dots\vertexgen_n$ is a path w.r.t.\ a digraph $\digraph$ iff every $(\vertexgen_i, \vertexgen_{i+1})\in\digraph$.
\emph{(In the formalisation, the path predicate and the vertex-set function are those of the library: \texttt{Vwalk.vwalk} / \texttt{vwalk\_bet} from \texttt{Directed\_Set\_Graphs.Vwalk} and \texttt{Pair\_Graph.dVs} from \texttt{Directed\_Set\_Graphs.Pair\_Graph}; the \texttt{DDFS.thy} source uses thin abbreviations \texttt{dpath}/\texttt{dpath\_bet} for backward compatibility.)}

A path $\vertexa\vertexb\dots\vertexgen_n$ is a simple path iff for every $1\leq i\neq j\leq n$, $\vertexgen_i \neq \vertexgen_j$.
For a vertex $\vertexgen$, $[\vertexgen]$ denotes a singleton path constituted of this vertex.
For two paths $\patha$ and $\pathb$, $\patha\cat\pathb$ denotes their concatenation.
If it causes no ambiguity, we refer to the set of vertices in a path with the path name.
E.g. for paths $\patha$ and $\pathb$, $\patha\cap\pathb$ refers to the intersection of the sets of vertices in the two paths.


In the following, we assume that we have an implicit fixed digraph $\DAG$, on which the algorithms operate.
We do not pass $\DAG$ as an argument to algorithms that do not change it for notational economy.
Lastly, we refer to the digraph $\DAG$ as a layered digraph iff it satisfies the following properties:
\begin{enumerate}
\item if $(\vertexa,\vertexb)\in\DAG$ then $\level(\vertexa) > \level(\vertexb)$, and
\item for any vertex $\vertexgen$ in the vertices of $\digraph$, there is a path $\path$ starting at $\vertexgen$ s.t. $\level(\last(\path)) = 0$.
\end{enumerate}
Any algorithms or theorems that state a digraph $\DAG$ in this section assume that it is a layered digraph.



\subsection{DFS}

\newcommand{\dfs}[0]{\textsf{\upshape dfs}}
\newcommand{\vis}[0]{\textsf{\upshape vis}}
\newcommand{\avoid}[0]{\textsf{\upshape avoid}}
\renewcommand{\invariant}[0]{\textsf{\upshape invariant}}
\newcommand{\none}[0]{\textsf{\upshape none}}

Let $\dfs$ be a function s.t. for a list of vertices $\path$, two sets of vertices $\avoid$ and $\vis$, and a predicate on vertices $P$; $\dfs(\avoid,\vis,\path,P)$ is a pair $(\path',\vis')$.
The first member of the pair, $\path'$, can be a list of vertices or nothing (denoted by $\none$), and the second member of the pair is a set of vertices.

\begin{myproperty}
  \label{property:dfs}
  For any $\avoid,\vis,\path$, and $P$, if $\dfs(\avoid,\vis,\path,P) = (\path',\vis')$, %% and $\invariant_0(\avoid,\vis,P)$ holds
the following statements hold:
  \begin{enumerate}
  \item {\label{property:dfs:1} $\invariant_0(\avoid,\vis'\setminus\vis,P)$}
  \item {\label{property:dfs:2}if $\path' = \none$, then $\path \setminus \avoid\subseteq \vis'$}
  \item if $\path'\neq\none$:
    \begin{enumerate}
      %\item for any path $\path''$ starting from a vertex $\vertexgen\in\path'$ and whose last vertex satisfies the predicate $P$, then $\path''\cap\avoid\neq\emptyset$
    \item {\label{property:dfs:12}$\path'$ is not empty}
    \item {\label{property:dfs:3}$\path'$ is a simple path w.r.t.\ $\DAG$}
    \item {\label{property:dfs:4}$P(\last(\path'))$ holds}
    \item {\label{property:dfs:5}$(\path'\setminus\butlast(\path))\cap(\avoid\cup\vis') = \emptyset$}
    \item {\label{property:dfs:6}$\head \path' = \head \path$}
    \item {\label{property:dfs:7}$\path\setminus\path'\subseteq\vis'$}
    \item {\label{property:dfs:11}if $\forall\vertexgen\in\butlast(\path).\;\neg P(\vertexgen)$, then we have $\forall\vertexgen\in\butlast(\path').\;\neg P(\vertexgen)$}
    \end{enumerate}
  \item {\label{property:dfs:8}$\vis \subseteq \vis'$}
  \item {\label{property:dfs:9}for any $\vertexgen\in\vis'\setminus\vis$, we have $\neg P (\vertexgen)$}
  \item {\label{property:dfs:10}for any $\vis'\setminus\vis\cap \avoid=\emptyset$}
  \item {\label{property:dfs:13}$\vis'\setminus \vis$ is finite}
  
  \end{enumerate}
\end{myproperty}

Above, $\invariant_0(\avoid,\vis,P)$ is a predicate that holds iff, for any vertex $\vertexgen\in\vis$ and a path $\path$ starting at $\vertexgen$, we have $\path\; \cap\; \avoid \neq \emptyset$ if $P(\last \path)$ holds, where $\last(\path)$ denotes the last vertex in a list of vertices $\path$.

Intuitively, $\dfs$ is a depth first search procedure that can start from a state of a previous call to itself, i.e. it can start from a stack of vertices currently being processed and a set of already visited vertices and continues searching from there.
It tries to find a path to a vertex that satisfies the given predicate, starting at a vertex in the given stack.
That path has to be disjoint from the set of vertices that is the first argument to $\dfs$.

\subsection{Amendments to the DFS interface}
\label{sec:dfs-amendments}

During the formal derivation of a $\dfs$ procedure satisfying Property~\ref{property:dfs}
(as an instance of a generic DFS skeleton; theories \textsf{DFS\_Avoid} and \textsf{IDFS}),
the property, as originally stated in the formalisation, turned out to be
\emph{unrealisable}: no function satisfies all of its clauses simultaneously, for any
digraph with at least one edge.  The following amendments repair this.  Throughout,
recall that in the formalisation the list $\path$ is kept in reverse order (its head is
the tip of the search stack), and that the digraph is assumed to be layered.

\begin{enumerate}
\item \textbf{A vestigial clause is dropped.}  The formalisation contained a clause with
  no counterpart in Property~\ref{property:dfs} (the fifth vis-works assumption), stating
  that every newly visited vertex $\vertexgen\in\vis'\setminus\vis$ satisfies
  $\vertexgen = \last(\path)$.  Together with clause~\ref{property:dfs:2}
  (completeness) this is contradictory: for $P = \lambda\vertexgen.\,\mathsf{False}$ the
  procedure must return $\none$, completeness then forces at least two distinct new
  visited vertices for a two-element stack, and both would have to equal $\last(\path)$.
  The clause was referenced only from a commented-out proof sketch and is deleted.

\item \textbf{The precondition on $\vis$ is reinstated.}  Clause~\ref{property:dfs:1}
  ($\invariant_0(\avoid,\vis'\setminus\vis,P)$) originally carried the precondition
  $\invariant_0(\avoid,\vis,P)$, which had been commented out.  Without it the clause is
  unrealisable in combination with clauses~\ref{property:dfs:5}, \ref{property:dfs:7}
  and~\ref{property:dfs:2}: those clauses force the search to prune at previously visited
  vertices and to record the pruned stack vertices as visited, so a $P$-vertex hidden
  behind $\vis$ (e.g.\ $\digraph\supseteq\{(\vertexa,\vertexb)\}$, $P=\{\vertexb\}$,
  $\vis=\{\vertexb\}$, $\avoid=\emptyset$) yields a newly visited vertex from which a
  $P$-vertex is reachable without meeting $\avoid$.  Clause~\ref{property:dfs:1} now
  carries the premises
  \[
    \invariant_0(\avoid,\vis,P), \qquad
    \forall\vertexgen\in\vis.\;\neg P(\vertexgen), \qquad
    \text{$\path$ is a simple path w.r.t.\ $\digraph$.}
  \]
  The path premise is needed because for an arbitrary (non-path) stack the newly visited
  vertices need not have been exhaustively explored; the layered structure of the digraph
  and the fact that the stack is a path exclude edges from finished vertices back into
  the live stack.

\item \textbf{Output simplicity needs input simplicity.}  Clause~\ref{property:dfs:3}
  ($\path'$ is a simple path) now carries the premise that $\path$ is simple: the
  procedure keeps the given stack as the bottom of its own stack, so duplicates in the
  input reappear in the output.

\item \textbf{Vertices of $\path$ in $\avoid$ are exempt from clause~\ref{property:dfs:7}.}
  The clause $\path\setminus\path'\subseteq\vis'$ becomes
  $(\path\setminus\path')\setminus\avoid\subseteq\vis'$.  A vertex of $\path$ that lies in
  $\avoid$ and is dropped from the stack can be recorded neither in $\vis'$ (it would
  violate clause~\ref{property:dfs:10}) nor kept in $\path'$
  (clause~\ref{property:dfs:5}).
\end{enumerate}

All premises added above are available at every call site of $\dfs$ inside
Algorithm~\ref{alg:ddfs}: invariants \ref{invar:2}~and~\ref{invar:3} are exactly the
$\invariant_0$ preconditions, invariants \ref{invar:10}~and~\ref{invar:11} give
$\neg P$ on the visited sets, and invariants \ref{invar:4}~and~\ref{invar:5} give that
the stacks are paths; simplicity of the stacks is maintained by
clause~\ref{property:dfs:3} itself together with the level structure.

\paragraph{Final form of the amendment to clause~\ref{property:dfs:1}.}
In the completed formalisation the history-style precondition of item~2 was replaced by
\emph{closure} clauses, which need no precondition at all: if
$\dfs(\avoid,\vis,\path,P) = (\none,\vis')$ then every edge from a vertex of
$\vis'\setminus\vis$ leads into $\avoid\cup\vis'$; if the result is
$(\path',\vis')$ then every such edge leads into $\avoid\cup\vis'\cup\path'$.
The $\invariant_0$-style facts are then derived by the caller by induction along walks.
Correspondingly, invariants \ref{invar:2}~and~\ref{invar:3} take a ``big'' form whose
hit-set also contains the searching side's own stack, and the termination measure of
Theorem~\ref{thm:ddfsterminates} is replaced by the lexicographic pair
$(\level(\vertexa)+\level(\vertexb),\;[\vertexa=\vertexb])$: the searches never
increase the level sum, and every search launched from a state with equal tips returns a
new tip outside its avoid set, hence unequal to the other tip.

A $\dfs$ procedure satisfying the amended property has been formally derived from a
generic DFS-skeleton locale: theory \textsf{DFS\_Avoid} proves the properties for the
skeleton instance, and theory \textsf{IDFS} transfers them to the exact vocabulary of
the DDFS development.  Under the amended interface the remaining unconditional clauses
are jointly realisable (degenerate inputs can always be answered by a $\none$ return, or
by a $\some$ return whose head is a $P$-vertex outside $\avoid\cup\vis$ when completeness
would otherwise force a $P$-vertex into the visited set).

\subsection{DDFS}

We now are ready to state the definition of the double depth first search (DDFS) procedure which operates on a layered digraph $\DAG$.
Assume there is a function $\level$ that maps every vertex in $\digraph$ to a natural number.
Also assume that the predicate $\Plevel{\vertexgen}$ is defined to be true for a vertex $\vertexgen'$ iff $\level(\vertexgen') \leq \level(\vertexgen)$.
DDFS takes as input two vertices from $\digraph$, and returns a pair of lists of vertices.
It is described as Algorithm~\ref{alg:ddfs}.
{
\SetKwIF{If}{ElseIf}{Else}{if}{}{else if}{else}{endif}
\begin{algorithm}\DontPrintSemicolon
  \begin{enumerate}
  \item $\patha=[]$; $\visa = \emptyset$; $\pathb=[]$; $\pathb'=[]$; $\visb = \emptyset$;\;
  \item \KwSty{while} $\level(\vertexa) \neq 0$ or $\level(\vertexb) \neq 0$\;
  \item \ \ \KwSty{if} $\vertexa = \vertexb$\;
  \item {\label{step:dfscall1}
        \ \ \ \ $(\path,\visb) = \dfs((\patha\cat[\vertexa])\cup\visa, \pathb'\cat[\vertexb], \visb, \Plevel{\vertexb})$\;}
  \item \ \ \ \ \KwSty{if} $\path \neq \none$\;
  \item {\label{step:newitr1}
        \ \ \ \ \ \ $\vertexb=\last(\path)$; $\pathb'=\butlast(\path)$\;}
  \item \ \ \ \ \KwSty{else}\;
  \item \ \ \ \ \ \ $\pathb=\pathb\cat\pathb'$; $\pathb'=[]$ //New barrier\;
  \item \ \ \ \ \ \ $(\path,\visa) = \dfs(\pathb\cat[\vertexb], \patha\cat[\vertexa],\visa, \Plevel{\vertexa})$\;
  \item \ \ \ \ \ \ \KwSty{if} $\path \neq \none$\;
        {\label{step:newitr2}
  \item \ \ \ \ \ \ \ \ $\vertexa=\last(\path)$; $\patha=\butlast(\path)$\;}
  \item \ \ \ \ \ \ \KwSty{else}\;
  \item {\label{step:ret1}
        \ \ \ \ \ \ \ \ \KwSty{return} $(\patha\cat[\vertexa], \pathb\cat[\vertexb])$}\;
  \item \ \ \KwSty{else} \KwSty{if} $\level(\vertexb) \leq \level(\vertexa)$\;
  \item {\label{step:newitr3}
        \ \ \ \ $\patha = \patha\cat[\vertexa]$; $\vertexa = \choice(\{\vertexgen\mid (\vertexa,\vertexgen)\in \digraph\})$\;}
  \item \ \ \KwSty{else}\;
  \item {\label{step:newitr4}
        \ \ \ \ $\pathb' = \pathb'\cat[\vertexb]$; $\vertexb = \choice(\{\vertexgen\mid (\vertexb,\vertexgen)\in \digraph\})$\;}
  \item {\label{step:ret2} \KwSty{return} $(\patha\cat[\vertexa], \pathb\cat\pathb'\cat[\vertexb])$}\;
  \end{enumerate}
  \caption{$\ddfs(\vertexa,\vertexb)$}\label{alg:ddfs}
\end{algorithm}
}
\newcommand{\patharet}{\ensuremath{\overline{\patha}}}
\newcommand{\pathbret}{\ensuremath{\overline{\pathb}}}

Before we state or prove the main theorems about DDFS, we first show that the following statements are invariants that hold throughout the execution of the while loop in Algorithm~\ref{alg:ddfs}.
\begin{enumerate}
\item {\label{invar:1} $\patha \cap \pathb\cat\pathb' = \emptyset$}
\item {\label{invar:2} $\forall\vertexgen\in\visa$ and $\path$, if $\head(\path) = \vertexgen$ and $\level(\last(\path)) \leq \level(\vertexb)$, then $\path\cap(\pathb\cat[\vertexb']) \neq \emptyset$, where $\vertexb'$ denotes $\head(\pathb'\cat[\vertexb])$}
\item {\label{invar:3} $\forall\vertexgen\in\visb$ and $\path$, if $\head(\path) = \vertexgen$ and $\level(\last(\path)) \leq \level(\vertexa)$, then $\path\cap((\patha\cat[\vertexa]) \cup \visa)\neq \emptyset$}
\item {\label{invar:4} $\patha\cat[\vertexa]$ is a path in the given digraph $\DAG$}
\item {\label{invar:5} $\pathb\cat\pathb'\cat[\vertexb]$ is a path in the given digraph $\DAG$}
\item {\label{invar:6} $\forall\vertexgen\in\patha. \level(\vertexa) < \level(\vertexgen)$}
\item {\label{invar:7} $\forall\vertexgen\in\pathb\cat\pathb'. \level(\vertexb) < \level(\vertexgen)$}
\item {\label{invar:6'} $\forall\vertexgen\in\patha$ either $\vertexb\neq\vertexgen$ or $\level(\vertexb) < \level(\vertexgen)$}
\item {\label{invar:7'} $\forall\vertexgen\in(\pathb\cat\pathb')$ either $\vertexa\neq\vertexgen$ or $\level(\vertexa) < \level(\vertexgen)$}
\item {\label{invar:8} $\pathb\setminus(\visa\cup(\patha\cat[\vertexa]))\subseteq\visb$}
\item {\label{invar:10} $\forall\vertexgen\in\visa.\level(\vertexgen)>\level(\vertexa)$}
\item {\label{invar:11} $\forall\vertexgen\in\visb.\level(\vertexgen)>\level(\vertexb)$}
\item {\label{invar:12} $\visa$ is finite}
\item {\label{invar:13} $\visb$ is finite}
%% \item \textcolor{red}{\label{invar:11} $\forall\vertexgen\in\visa.\level(\vertexgen)>\level(\vertexa)$}
%% \item \textcolor{red}{\label{invar:12} $\forall\vertexgen\in\visb.\level(\vertexgen)>\level(\vertexb)$}
%% \item \textcolor{red}{\label{invar:13} $(\pathb\cat\pathb')\cap\visa=\emptyset$}
%% \item \textcolor{red}{\label{invar:13'} if $\pathb\cat\pathb'$ is empty, then $\vertexb\not\in\visa$}
\end{enumerate}

\begin{mythm}
  \label{thm:ddfsterminates}
  For any two vertices $\vertexa,\vertexb\in\vertices(\DAG)$, $\ddfs(\vertexa,\vertexb)$ terminates.
\end{mythm}
%% Before we prove any properties about the lists returned by DDFS, it should be clear that the termination of Algorithm~\ref{alg:ddfs} follows from invariants~\ref{invar:8},~\ref{invar:10},~and~\ref{invar:11}.
\begin{proof}
To show termination, we show that the quantity $\level(\vertexa)+\level(\vertexb)+\cardinality{\vertices(\DAG)\setminus\visa}+\cardinality{\vertices(\DAG)\setminus\visb}$ decreases with every iteration of the algorithm.
The algorithm starts new iterations at four points: at the end of lines~\ref{step:newitr1},~\ref{step:newitr2},~\ref{step:newitr3},~and~\ref{step:newitr4}.
We shall analyse each of them separately.
Note: if a variable $x$ is assigned a new value in an iteration, we refer to it as $x^*$ after being assigned the new value.
This is only during the iteration in which the new value was assigned.
E.g. after line~\ref{step:dfscall1}, we refer to the set $\visb$ as $\visb^*$ in our proofs.

{\bf Line~\ref{step:newitr1}:} from invariant~\ref{invar:11}, we have that $\vertexb\not\in\visb$ before the iteration starts.
However, from properties~\ref{property:dfs:5}~and~\ref{property:dfs:7}, we have that $\vertexb\in\visb^*$ after the call to $\dfs$ on line~\ref{step:dfscall1}, since $\vertexb$ was in the avoid argument to $\dfs$.
Accordingly, the size of the set $\visb^*$ is at least one bigger than that of $\varb$, i.e. its size at the beginning of the iteration.
This implies that the expression above is one less in the new iteration.

{\bf Line~\ref{step:newitr2}:} this uses an argument analogous to that used for the case when the new iteration is started on line~\ref{step:newitr1}, but regarding the size of the set $\visa$ instead of $\visb$.

{\bf Line~\ref{step:newitr3}:} firstly, there will always be an edge from $\vertexa$ to another vertex, since $\DAG$ is a layered digraph and since the layer of $\vertexa$ is not $0$, which comes from the condition of the while loop.
Also, we have that $\level(\vertexa^*) < \level(\vertexa)$, since $(\vertexa, \vertexa^*)\in\DAG$ and $\DAG$ is a layered digraph.
Thus, the expression above will be less in the next iteration.

{\bf Line~\ref{step:newitr4}:} this is analogous to the case of line~\ref{step:newitr3}.
\end{proof}





\begin{mythm}
  \label{thm:ddfsvalid}
  Assume that $\dfs$ satisfies Property~\ref{property:dfs}.
  %, for any two sets of vertices $\avoid$ and $\vis$, path $\path$ and a predicate $P$ on vertices from the digraph $\DAG$.

  Then, for any two vertices $\vertexa$ and $\vertexb$ in the digraph $\digraph$, if $\ddfs(\vertexa,\vertexb) = (\patharet,\pathbret)$, then one of the following statements holds:
  \begin{enumerate}
  \item $\last(\patharet) = \last (\pathbret)$ and $\butlast(\patharet) \cap \butlast(\pathbret) = \emptyset$ and $\forall\path.$ if $\path$ is a path starting at $\vertexa$ or $\vertexb$ s.t. $\level(\last(\path)) \leq \level(\last(\patharet))$, then $\last(\patharet)\in\path$
  \item $\patharet \cap \pathbret = \emptyset$ and $\level(\last(\patharet)) = \level(\last(\pathbret)) = 0$
  \end{enumerate}
  Furthermore, $\patharet$ and $\pathbret$ are both paths in the given digraph $\DAG$.
\end{mythm}
To prove this theorem, we need the following lemmas.

\begin{mylem}
\label{lem:ddfs:1}
At line~\ref{step:ret1}, for any path $\path$ starting at a vertex $\vertexgen\in(\pathb\cat[\vertexb])$ s.t.\ $\level(\last(\path)) \leq \level(\vertexa)$, there are $\path',\path''$, and $\vertexgen'$ s.t.\ $\path=[\vertexgen]\cat\path'\cat[\vertexgen']\cat\path''$ and $\vertexb\in(\pathb\cat[\vertexb])$.
\end{mylem}
\begin{proof}
From our assumption on $\path$, and invariants~\ref{invar:3}~and~\ref{invar:8}, we have that $\path\cap(\patha\cup\visa^*)\neq\emptyset$.
Also, from Property~\ref{property:dfs:2} of $\dfs$, we have that $\patha\subseteq\visa^*$ since $\dfs$, when called at step~\ref{step:dfscall1}, was unable to find a path starting at a vertex on $\patha\cat[\vertexa]$ and from $\dfs$.
Accordingly, there exists two paths $\path_a$ and $\path_b$ and a vertex $\vertexgen_a\in\visa^*$ s.t.\ $\path=\path_a\cat[\vertexgen_a]\cat\path_b$.
From that and from invariant~\ref{invar:2}, we have that $[\vertexgen_0]\cat\path_b\cap(\patha\cat[\vertexb]) \neq \emptyset$.
Thus, there are paths $\path_c$, $\path_d$, and $\vertexgen_b$ s.t.\ $\path_b=\path_c\cat[\vertexgen_b]\cat\path_d$ and $\vertexgen_b\in(\pathb\cat[\vertexb])$, which shows that $\path=[\vertexgen]\cat\path'\cat[\vertexgen']\cat\path''$, for some $\vertexgen'\in(\pathb\cat[\vertexb])$, if we set $\path'=\path_a\cat[\vertexgen_a]\cat\path_c$, $\vertexgen'=\vertexgen_b$, and $\path''=\path_d$.
\end{proof}

\begin{mylem}
\label{lem:ddfs:2}
For a predicate $P$ and a list $\path \neq []$, suppose we have $\path_b'$, $\vertexgen'$ and $\path_b''$, s.t. $\path_b=\path_b'\cat[\vertexgen']\cat\path_b''$ and $P(\vertexgen')$ hold, for any $\path_a,\vertexgen,\path_b$, if $\path=\path_a\cat[\vertexgen]\cat\path_b$, $\path_b\neq[]$ and $P(\vertexgen)$.
Then $P(\last(\path))$ holds.
\end{mylem}

\begin{mylem}
\label{lem:ddfs:3}
At line~\ref{step:ret1}, for any path $\path$ starting at a vertex $\vertexgen\in(\pathb\cat[\vertexb])$, then $\last(\path)\in(\pathb\cat[\vertexb])$.
\end{mylem}
\begin{proof}
This follows from Lemma~\ref{lem:ddfs:1}~and~\ref{lem:ddfs:2}.
\end{proof}


\begin{mylem}
\label{lem:ddfs:4}
At line~\ref{step:ret1}, for any path $\path$ starting at a vertex $\vertexgen\in(\pathb\cat[\vertexb])$ s.t.\ $\level(\last(\path)) \leq \level(\vertexa)$, we have that $\vertexb\in\path$.
\end{mylem}
\begin{proof}
From Lemma~\ref{lem:ddfs:3}, we know that $\last(\path) \in (\pathb\cat[\vertexb])$.
Also, observe that our assumptions on $\DAG$ imply that for any vertex $\vertexgen\in\butlast(\path)$, $\level(\vertexgen) > \level (\last (\path))$ holds.
Accordingly, the fact that $\level(\last(\path)) \leq \level(\vertexa)$ implies that $\vertexb\in\path$.
\end{proof}


\begin{mylem}
\label{lem:ddfs:5}
At line~\ref{step:ret1}, for any path $\path$ starting at a vertex $\vertexgen\in(\patha\cat[\vertexa])$ s.t.\ $\level(\last(\path)) \leq \level(\vertexa)$, we have that $\vertexb\in\path$.
\end{mylem}
\begin{proof}
From Property~\ref{property:dfs:2} of $\dfs$, and the fact that $\dfs$ could not replace $\patha\cat[\vertexa]$ when , we have that $\patha\cat[\vertexa]\subseteq\visa$.
Accordingly, $\vertexa\in\visa$.
From that, from invariant~\ref{invar:2}, and the fact that $\pathb'=[]$, we have that $\path\cap(\pathb\cat[\vertexb])\neq\emptyset$.
Hence, there is a vertex $\vertexgen'\in(\pathb\cat[\vertexb])$; and paths $\path'$ and $\path''$ s.t. $\path=[\vertexgen]\cat\path'\cat[\vertexgen']\cat\path''$.
From that and from Lemma~\ref{lem:ddfs:4}, we have that $\vertexb\in\path$.
\end{proof}

\begin{mylem}
\label{lem:ddfs:6}
$\patha\cap(\pathb\cat\pathb') = \emptyset$
\end{mylem}
\begin{proof}
This follows from invariant~\ref{invar:1}.
\end{proof}

\begin{proof}[Proof of Theorem~\ref{thm:ddfsvalid}]
The first property of DDFS we prove is that both $\patharet$ and $\pathbret$ are paths in $\DAG$.
This follows from invariants~\ref{invar:4}~and~\ref{invar:5}.

To prove the other statement about the returned paths from Algorithm~\ref{alg:ddfs}, we consider the two cases in which it returns separately.

{\bf Return at step~\ref{step:ret1}}: from Lemmas~\ref{lem:ddfs:4},~\ref{lem:ddfs:5},~and~\ref{lem:ddfs:6}, and the fact that $\vertexa=\vertexb$, we have that
$\last(\patharet) = \last (\pathbret)$ and $\butlast(\patharet) \cap \butlast(\pathbret) = \emptyset$ and $\forall\path.$ if $\path$ is a path starting at $\vertexa$ or $\vertexb$ s.t. $\level(\last(\path)) \leq \level(\last(\patharet))$, then $\last(\patharet)\in\path$.
This is what we want to prove for the paths returned by DDFS.

{\bf Return at step~\ref{step:ret2}:} from invariant~\ref{invar:1}, we have that $\patha\cap(\pathb\cat\pathb')=\emptyset$.
From that, and the fact that $\vertexa\neq\vertexb$, we have that $\patharet$ and $\pathbret$, are disjoint.

\end{proof}

\subsection{Formalisation details}

The following is the recursive definition of DDFS that is actually used to prove theorems about DDFS by computational induction.

\SetKwIF{If}{ElseIf}{Else}{if}{}{else if}{else}{endif}
\begin{algorithm}\DontPrintSemicolon
  \KwSty{if} $\vertexa = \vertexb$\;
  \ \ $(\path,\visb) = \dfs(\patha\cat[\vertexa], \pathb'\cat[\vertexb], \visb, \Plevel{\vertexb})$\;
  \ \ \KwSty{if} $\path \neq \none$\;
  \ \ \ \ $\vertexb=\last(\path)$\;
  \ \ \ \ $\pathb'=\butlast(\path)$\;
  \ \ \ \ \KwSty{return} $\ddfs(\vertexa,\path_1,\visa,\vertexb,\path_2,\path_2',\visb)$\;
  \ \ \KwSty{else}\;
  \ \ \ \ $\pathb=\pathb\cat\pathb'$ //New barrier\;
  \ \ \ \ $\pathb'=[]$\;
  \ \ \ \ $(\path,\visa) = \dfs(\pathb\cat[\vertexb], \patha\cat[\vertexa],\visb, \Plevel{\vertexa})$\;
  \ \ \ \ \KwSty{if} $\path \neq \none$\;
  \ \ \ \ \ \ $\vertexa=\last(\path)$\;
  \ \ \ \ \ \ $\patha=\butlast(\path)$\;
  \ \ \ \ \ \ \KwSty{return} $\ddfs(\vertexa,\path_1,\visa,\vertexb,\path_2,\path_2',\visb)$\;
  \ \ \ \ \KwSty{else}\;
  \ \ \ \ \ \ \KwSty{return} $(\patha\cat[\vertexa], \pathb\cat[\vertexb])$\;
  \KwSty{else} \KwSty{if} $\level(\vertexa) = 0$ and $\level(\vertexb) = 0$\;
  \ \ \KwSty{return} $(\patha\cat[\vertexa], \pathb\cat\pathb'\cat[\vertexb])$\;
  \KwSty{else} \KwSty{if} $\level(\vertexb) \leq \level(\vertexa)$\;
  \ \ $\patha = \patha\cat[\vertexa]$\;
  \ \ $\vertexa = \choice(\{\vertexgen\mid (\vertexa,\vertexgen)\in \digraph\})$\;
  \ \ \KwSty{return} $\ddfs(\vertexa,\path_1,\visa,\vertexb,\path_2,\path_2',\visb)$\;
  \KwSty{else}\;
  \ \ $\pathb' = \pathb'\cat[\vertexb]$\;
  \ \ $\vertexb = \choice(\{\vertexgen\mid (\vertexb,\vertexgen)\in \digraph\})$\;
  \ \ \KwSty{return} $\ddfs(\vertexa,\path_1,\visa,\vertexb,\path_2,\path_2',\visb)$\;
  \caption{$\ddfs(\vertexa,\path_1,\visa,\vertexb,\path_2,\path_2',\visb)$}\label{alg:BlossomOrAug}
\end{algorithm}
